\chapter{Summary}\label{ch:work-brief}

This document describes the instruction encoding and semantics of the Apple M5 GPU (the G17) and the path by which
work reaches it below Metal. From that description the project built a compiler that emits native
G17 code, and two ways to run it: one through Metal, one below it, where the project's runtime submits
work through Apple's private IOGPU interfaces and Metal is never loaded (\cref{fig:bypass}). The measurements
come from one Apple M5 Pro with 20 GPU cores and 48~GB of memory (MM~\mm{22}, MM~\mm{25.91}); \cref{sec:measured-part-software}
gives the platform and which records name the OS build.\label{sec:part-software}

\section{Recovered behavior}\label{sec:has-been-recovered}

Apple documents the GPU and its tensor unit, the neural accelerator, only at the Metal interface. Below that
interface:

\begin{itemize}
\item The usable register namespace ends at R125, and an instruction naming any register from R126 to R143
  is dropped whole, without an error (\cref{sec:namespace}). Register operands carry a lifetime modifier, and in a few
  forms the encoding fixes it (\cref{tab:fixed-modifiers}). A \emph{release} clears the lane's value as it is read.
  Controlled probes then read zero, but that is not established in general, so a later reader must not rely on it
  (\cref{sec:release-does}).
\item Instruction lengths are even, from 2 to 22 bytes. Apple's code uses every even length from 2 to 16, plus 18- and
  22-byte forms such as \op{10909} and \op{14661}; no 20-byte instruction has been seen. One opcode can appear at several
  lengths. Operands are linear functions of
  scattered bits, recovered as field maps. The tensor-instruction encoder reproduces 800 of 800 compiler-emitted
  instructions byte for byte. Of the \val{named-instructions} opcodes the project names, \val{executed-instructions} are confirmed by execution
  (\cref{ch:instruction-encoding}).
\item Apple's decoder admits ten tensor MMA forms. The mapping from each lane's registers to matrix elements
  was measured with one-hot probes over 6,651 dispatches. For half inputs, the exact reduction order of one 16 × 16 × 16
  multiply-accumulate, including where the accumulator enters, reproduces 5,120 of 5,120 tile elements bit for bit.
  A sequential chain, a balanced tree and an exactly rounded sum each match fewer than 2,800. For bfloat and
  truncated-fp32 inputs only a few extreme cases were measured (\cref{sec:arithmetic-semantics-supported}). Clearing one
  metadata byte, the unit's enable, makes every tensor MMA write zeros (\cref{sec:ordering-enable}). Ordering uses no
  wait instruction. The producer of a late value names one of eight scoreboard slots, and the consumer carries an
  eight-bit wait mask; if the slot's bit is left out, the consumer reads the register before the value lands, and in
  the measured tensor case it reads zeros (\cref{ch:synchronization}).
\item Every byte across 421 delivered archives has a known source. Apple's loader requires 148
  constants, and the meaning of 12 of their bytes is still open. Metal checks the container's format and does not inspect the code in it (\cref{ch:kernel-metadata-image}).
\item The launch words, the Submit record, the kernel entry and the completion path below Metal are recovered
  (\cref{ch:submission-below-metal}). With them the native runtime runs MiniLM-L6 retrieval and Qwen2.5-0.5B generation
  without Metal, on resource graphs of the four measured sizes, from 3~MiB to 1~GiB (\cref{sec:measured-graphs-generalizes}).
\end{itemize}

\section{Two ways to run a program}\label{sec:project-authors-two}

The project builds the object, metadata and container around the native instruction bytes, an executor
for each path (\code{tools/g17decodegen.m} through Metal, the native runtime below it) and the checkers, among them the
released-register check and the CPU emulator (\cref{app:d}).

% The two execution paths. Blue: the project's components. Red: the
% project's native runtime. Grey: Apple's components.
\begin{figure}[htbp]
\centering
\begin{tikzpicture}[node distance=6mm and 8mm]
  \node[ours, minimum width=46mm] (ir) {IR and kernel builders};
  \node[ours, minimum width=46mm, below=of ir] (cc) {Compiler};
  \node[ours, minimum width=46mm, below=of cc] (code) {Native G17 instruction bytes};
  \draw[flow] (ir) -- (cc);
  \draw[flow] (cc) -- (code);

  \node[ours, text width=52mm, below left=9mm and -14mm of code] (img)
    {\mbox{Metal-compatible} image\\[-1pt]{\footnotesize object, metadata, library, archive}};
  \node[ours, text width=52mm, below=of img] (exec)
    {Metal-path executor\\[-1pt]{\footnotesize \code{decodegen}}};
  \node[apple, text width=52mm, below=of exec] (metal)
    {Apple Metal\\[-1pt]{\footnotesize pipeline creation and submission}};

  \node[nat, text width=52mm, below right=9mm and -14mm of code] (rt)
    {Native runtime\\[-1pt]{\footnotesize allocation, resource and launch state,}\\[-2pt]
     {\footnotesize submission records, completion handling}};

  \draw[flow] (code.south) -- ++(0,-4mm) -| (img.north);
  \draw[flow] (code.south) -- ++(0,-4mm) -| (rt.north);
  \draw[flow] (img) -- (exec);
  \draw[flow] (exec) -- (metal);

  \path (metal.south) ++(0,-9mm) coordinate (iogpuTop);
  \node[band, below=9mm of metal.south -| code, minimum width=118mm] (iogpu)
    {Apple private IOGPU framework and IOKit interfaces};
  \node[band, below=5mm of iogpu, minimum width=118mm] (drv) {Apple AGX kernel driver};
  \node[band, below=5mm of drv, minimum width=118mm] (fw) {Apple GPU firmware};
  \node[box, draw=ink, fill=white, below=5mm of fw, minimum width=118mm] (gpu)
    {M5 GPU: scalar and tensor execution};

  \draw[flow] (metal.south) -- (metal.south |- iogpu.north);
  \draw[->, line width=1.4pt, color=native] (rt.south) -- (rt.south |- iogpu.north)
    node[pos=0.13, right=1mm, tag, text=native, text width=26mm, align=left]
      {direct submission};
  \draw[flow] (iogpu) -- (drv);
  \draw[flow] (drv) -- (fw);
  \draw[flow] (fw) -- (gpu);

  % Path labels
  \node[font=\small\bfseries, text=ours, anchor=base west, inner sep=1pt] at ([yshift=2mm]img.north west) {Through Metal};
  \node[font=\small\bfseries, text=native, anchor=base east, inner sep=1pt] at ([yshift=2mm]rt.north east) {Below Metal};

  % The Metal layer, which the native path does not enter
  \begin{scope}[on background layer]
    \node[draw=apple, dashed, rounded corners=3pt, fill=apple!4, inner sep=3mm,
          fit=(metal) (metal -| rt.east)] (mlayer) {};
  \end{scope}
  \node[tag, text width=44mm] at ($(metal.east)!0.5!(metal.east -| rt.south)$) {Metal not loaded in\\the native runtime};

  % Brace over the shared Apple layers
  \draw[decorate, decoration={brace, amplitude=5pt}, line width=0.6pt, color=apple]
    ([xshift=3mm]iogpu.north east) -- ([xshift=3mm]fw.south east)
    node[midway, right=6pt, tag, text width=18mm, align=left] {below both\\paths};
\end{tikzpicture}
\caption{The two execution paths. The native runtime does its own pipeline setup and submits directly through IOGPU.
Grey: Apple's; blue: the project's; red: the project's native runtime.}
\label{fig:bypass}
\end{figure}


Through Metal (\cref{ch:through-metal}), the project hands Metal an image, and Metal builds the pipeline from its
archived code and submits the work. Below Metal (\cref{ch:below-metal-native}), the native runtime takes the program
bytes directly. It uploads them, writes each launch's resource and launch state, builds the Submit records and calls
IOGPU itself; no Metal pipeline, command buffer or device exists in that process. Both paths run on IOGPU, the AGX
kernel driver and the firmware.

The compiler's encoders, its whole-program
checks and the CPU emulator decode programs with Apple's G17 decoder, loaded from \code{GPUCompiler.framework}
(\cref{ch:instruction-encoding}). That decoder is not used when the GPU runs a program, but it is the reason compiling and
validating currently require macOS. Metal source first passes through Apple's \code{xcrun metal} front end on its way to the
project's IR (\cref{ch:source-native-program}).

\section{The running example}\label{sec:one-program-ir}

The running example computes a 17 × 19 × 16 matrix product, half-precision inputs and fp32 output, and then adds 1.0
to one output element. Recompiling the IR reproduces the program that ran below Metal byte for byte: 1,232 bytes, SHA-256 \hash{8bb67a3c65919cb7fa9a6f2e70d152cb2d8086fba0cade7437f18d1d8114c667},
register count 76.

The source is seven lines of the project's IR: three buffers, \code{tensor\_matmul(A, B, C, M=17, N=19, K=16)}, a
load, an fp32 add and a store (\cref{ch:running-example-through}).

The 17 × 19 output does not divide into 16 × 16 tiles, so the product is two tile rows by two tile columns, and
the second row and column are partial (\cref{fig:example-tiles}).

% The running example's 17 x 19 output against the 16 x 16 tiles that cover it.
\begin{figure}[htbp]
\centering
\begin{tikzpicture}[x=1mm, y=1mm,
    lbl/.style={tag, font=\scriptsize, align=center},
    num/.style={circle, draw=ink!70, fill=white, inner sep=0.8pt, font=\scriptsize, text=ink},
    key/.style={font=\scriptsize, text=ink, anchor=west}]

  \def\u{1.9}   % one element
  \def\g{60.8}  % 32 elements: the two tile rows and columns together

  % The part of the tile grid the matrix actually reaches.
  \fill[ours!30] (0,0) rectangle ({19*\u},{-17*\u});

  % The whole 32 x 32 tile grid, and the tile boundaries inside it.
  \draw[ink!35, line width=0.4pt, dashed] (0,0) rectangle ({\g},{-\g});
  \draw[ink!35, line width=0.4pt, dashed] ({16*\u},0) -- ({16*\u},{-\g});
  \draw[ink!35, line width=0.4pt, dashed] (0,{-16*\u}) -- ({\g},{-16*\u});

  % The live result on top of it.
  \draw[ink!75, line width=0.7pt] (0,0) rectangle ({19*\u},{-17*\u});
  \draw[ink!75, line width=0.7pt] ({16*\u},0) -- ({16*\u},{-17*\u});
  \draw[ink!75, line width=0.7pt] (0,{-16*\u}) -- ({19*\u},{-16*\u});

  \node[num] at ({8*\u},{-8*\u})  {1};
  \node[num] at ({24*\u},{-8*\u}) {2};
  \node[num] at ({8*\u},{-24*\u}) {3};
  \node[num] at ({24*\u},{-24*\u}){4};

  \draw[decorate, decoration={brace, amplitude=4pt}, line width=0.5pt, color=ink!70]
    (0,1.5) -- ({19*\u},1.5) node[midway, above=4pt, lbl] {19 columns};
  \draw[decorate, decoration={brace, amplitude=4pt, mirror}, line width=0.5pt, color=ink!70]
    (-1.5,0) -- (-1.5,{-17*\u}) node[midway, left=4pt, lbl] {17\\rows};

  % Key, clear of the drawing.
  \node[key] at ({\g+9},-4)  {\textbf{1}\ \ 16 × 16, the only full tile\ \ R24--R31};
  \node[key] at ({\g+9},-11) {\textbf{2}\ \ 16 × 3\ \ R32--R39};
  \node[key] at ({\g+9},-18) {\textbf{3}\ \ 1 × 16\ \ R40--R47};
  \node[key] at ({\g+9},-25) {\textbf{4}\ \ 1 × 3\ \ R48--R55};
  \node[key, text width=52mm, align=left, text=ink!65] at ({\g+9},-36)
    {Shaded: the result. Dashed: the rest of each 16 × 16 tile, which the 20 mask instructions
     switch off.};
\end{tikzpicture}
\caption{The running example's output against the tiles that cover it. A 17 × 19 result does not divide into
16 × 16, so it takes two tile rows by two tile columns, three of them partial: one \op{5107} and one
eight-register accumulator group each.}
\label{fig:example-tiles}
\end{figure}

 The compiler emits 101 instructions, grouped by phase in \cref{tab:example-phases}.

\begin{xltabular}{\linewidth}{@{}l>{\hsize=1.10\hsize}X>{\hsize=0.90\hsize}X@{}}
\caption{The example program's 101 instructions by phase, as read with Apple's decoder.}\label{tab:example-phases}\\
\toprule
Phase & Instructions & What they do \\
\midrule
\endfirsthead
\tablecontinued{3}
\toprule
Phase & Instructions & What they do \\
\midrule
\endhead
addresses & one \op{14060} reading the lane index, then
shifts, ANDs and adds & compute each lane's row, column and byte offset
(\cref{sec:fragment-layouts}) \\*
masks & 20 \op{612} and \op{437} & enable only the 17 rows and 19 columns that exist
(\cref{sec:masks-predicated-memory}) \\*
loads & 4 \op{12674} and 4 \op{12675} & gather two A and two B fragments, eight bytes per lane each \\*
multiply & 4 \op{5107} & one per output tile, since K = 16 is one slice \\*
stores & 2 \op{17257} and 6 \op{17258} & write the four tiles, masked where partial \\*
epilogue & \op{14059}, \op{12682}, \op{11842}, \op{998}, \op{17229}, \op{684} & add 1.0 to one element \\*
\bottomrule
\end{xltabular}

The first \op{5107} is the ten bytes \hash{a70405112204a4620202}, which Apple's decoder prints as:

\begin{lstlisting}
op5107  R24..R31 <- 0x01000000, 1, R20_R21_R22_R23, 0, 2, R68_R69_R70_R71, 0, 2
        D        operand 1         A fragment      lifetime, type   B fragment   lifetime, type
\end{lstlisting}

\begin{itemize}
\item \textbf{D, R24–R31,} is an eight-register fp32 accumulator group holding the output tile at application rows
  0–15, columns 0–15. The compiler packs every tile with the hardware's B lane map, so its labeling of these registers
  differs from the instruction's canonical one by a one-bit rotation of the row index. A is packed the same way, and
  the product is correct in application coordinates (\cref{sec:coordinate-systems}).
\item \textbf{A, R20–R23, and B, R68–R71,} are four-register fragments of 16-bit values. Type code 2 means half.
\item \textbf{Operand 1, 0x01000000,} is the destination's raw operand, which older notes call the flags word. Its
  bit 24 means wait for scoreboard slot 0. All eight tensor loads write slot 0. Each
  load's own control word, 0x01100000, both publishes slot 0 and waits on it (\cref{sec:scoreboard-publish-wait,sec:scoreboard}). On an MMA
  whose loads fill slot 7, clearing the matching wait bit made the multiply read its fragments before the loads landed,
  and it produced zeros (measured on a separate slot-7 program, MM~\mm{6}).
\item \textbf{The 1 that follows it} is constant; no encoding bit varies it (\cref{sec:field-maps}).
\end{itemize}

Each fragment is read by two of the four MMAs, and the second read releases it (\cref{tab:example-mma-releases}). Each output element sums 16
products in a fixed tree, adjacent pairs and then pairs eight apart, rounding to nearest-even fp32 at every step (the
rule fitted on half inputs, which this program uses). The program uses the from-zero form, with no C input, so the
sum starts from the first of the tree's four partial sums; the accumulating form adds C first
(\cref{sec:arithmetic-semantics-supported}).

The epilogue orders its late values through the scoreboard.

\begin{enumerate}
\def\labelenumi{\arabic{enumi}.}
\item The \op{14059} of the threadgroup position publishes slot 0.
\item The \op{12682} waits on slot 0 for its index register and publishes slot 7.
\item The \op{998} waits on slot 7 for the loaded value.
\item The store needs no wait: its value comes from an ALU instruction, and ALU results are interlocked (\cref{sec:dependences-issue-scalar}).
\end{enumerate}

Alongside the code, the compiler reports what an executor needs: the three bindings with their element types, which
binding is written, the 24 distinct (opcode, length) forms used, an empty constant pool, the register count and
the system registers read. The project's image-building code turns that into the object, metadata, library and archive Metal loads (\cref{ch:kernel-metadata-image}).

The program was then run on the GPU below Metal (measured; \code{docs/g17-first-dispatch-trials.md}, "A second authored tensor program
below Metal"). The native runtime wrote the program at allocation 0, offset 0x6c0. Relative to the first tensor
program it had run, the resource graph changed in nine fields: two view offsets, one C pointer, two binding pointers, the kernel
word at offset 0x3f8 (from 0xa to 0xc), and the A, B and C windows, which became 544, 608 and
1,292 bytes, exactly 17 × 16 halves, 16 × 19 halves and 17 × 19 floats.

Three fresh processes each submitted the program through IOGPU with no Metal loaded. In each, 323 of 323 fp32 outputs
were bit-exact, inputs and guards were intact, both completion callbacks arrived and no recovery event occurred. The
outputs hash to \hash{ec43027260f5a0b8dd59526c17102ffe00066852f10af42a91103c36e778faf6}, the same as an independently
computed reference and as the earlier run through Metal (\cref{ch:submission-below-metal}).

\section{Applications}\label{sec:understanding-supports}

InternLM2.5-1.8B decodes on the compiler's kernels, run through Metal. Single stream at a 1,024-token context, it runs
at 1.14× the decode rate of mlx-lm's \code{generate\_step}; the
\href{https://github.com/sbryngelson/AGXForge/blob/main/docs/demonstrations.md}{demonstrations document} gives the other models and
workloads with their baselines. MiniLM-L6 sentence retrieval and Qwen2.5-0.5B-Instruct generation ran through the native
submission path. Each stage was checked against a CPU recomputation from its GPU inputs, and each output against the
original checkpoint evaluated in FP64 (MM~\mm{25.210}; \cref{ch:below-metal-native}). The compiler
was also run on code it was not written for. Of the 125 sources that compiled from a fixed set of 197 (the census) when it
was run on the hardware, 101 have run and been checked, and none crashes the driver (MM~\mm{25.149.1}); today's front end
and backend compile 128 of the 197 (\cref{ch:source-native-program}).

A CPU emulator, working from the program bytes, executes one position of a shipped batched-decode graph and of
single-stream decode, dispatch by dispatch. All 316 and all 172 dispatches match the GPU byte for byte, as does a
16~MB prefill slice (MM~\mm{25.197}).

\section{Limits of the evidence}\label{sec:where-evidence-stops}

The InternLM2 decode lead over mlx-lm (\cref{sec:understanding-supports}) comes from the kernels. In a matched study,
Apple's compiler, given the same kernels, ran them 10–15 percent faster with identical tokens, so the lead does not show
that native instruction control is needed. The tensor-unit kernels were not part of that study (MM~\mm{25.211}).

Below Metal, MiniLM and Qwen run with checked results, but more slowly than Apple's paths: the matched MiniLM FFN and
attention blocks take 2.427 and 3.815~ms against Metal's 1.797 and 1.962~ms (MM~\mm{25.210.5}). One auxiliary control
in the native-path validation fails. Qwen's prefix logits, compared with an FP64 model that uses the same half-precision
projections, deviate by up to 0.097, and 10,658 elements fall outside the bound \code{0.025 + 0.002 * abs(reference)}
(\cref{ch:below-metal-native}).

Because validation uses Apple's decoder (\cref{sec:project-authors-two}), and the project's own decoder for its
programs was started and then paused, nothing here establishes portability (MM~\mm{25.209}).

\section{Conventions}\label{sec:read-this-document}

Rules and measurements in Part~\ref{part:i} carry one of these evidence tags, which say where a statement stands in
\cref{fig:evidence-flow}.

% Where the rules in this document come from, and what they constrain.
% (Mermaid 0.13 of the machine model, drawn properly.)
\begin{figure}[htbp]
\centering
\begin{tikzpicture}[node distance=5mm and 11mm,
    src/.style={box, draw=ink!60, fill=ink!5, text width=38mm, align=left, font=\scriptsize,
      minimum height=15mm},
    pol/.style={box, draw=ours, fill=ours!10, text width=38mm, align=left, font=\scriptsize,
      minimum height=15mm},
    run/.style={box, draw=native, fill=native!10, text width=38mm, align=left, font=\scriptsize,
      minimum height=13mm}]

  \node[src] (hw)
    {\textbf{Measured hardware facts}\\[-1pt]opcodes, register reach, layouts, arithmetic, the
     scoreboard, throughput};
  \node[src, below=of hw] (ap)
    {\textbf{Apple's compiled code}\\[-1pt]AIR lowering, spills, library schedules, fast-math
     reassociation};

  \node[pol, right=of hw] (cc)
    {\textbf{This compiler's policy}\\[-1pt]admitted shapes, the memory bridge, named refusals,
     reserved registers};
  \node[pol, right=of ap] (abi)
    {\textbf{Measured Metal object ABI}\\[-1pt]slot 44's enable, slot 29's system-register
     vector, bindings, entry and prologue};

  \node[run, right=of cc, yshift=-11mm] (rt)
    {\textbf{Runtime policy}\\[-1pt]one simdgroup, exact grid, guards and reference checks};

  \draw[flow] (hw) -- (cc);
  \draw[flow] (ap) -- (cc);
  \draw[flow] (hw) -- (abi);
  \draw[flow] (ap) -- (abi);
  \draw[flow] (cc) -- (rt);
  \draw[flow] (abi) -- (rt);

  \node[tag, font=\scriptsize, anchor=north west, text width=82mm] at ([yshift=-7mm]ap.south west)
    {Grey: what the machine and Apple's code were observed to do. Blue: what this project then
     chose. A rule's tag says which of the two it rests on, and \emph{open} marks a question
     neither settles.};
\end{tikzpicture}
\caption{Where the rules in Part~\ref{part:i} come from. Measurement and observation of Apple's code constrain the
compiler's policy and the object ABI, and those two constrain what the runtime may do; the evidence tags above name
which kind of standing each rule has.}
\label{fig:evidence-flow}
\end{figure}


\begin{itemize}
\item \emph{measured}: run on this GPU and compared with an independent reference, with a control that could have failed;
\item \emph{measured (receipt fit)}: a recorded GPU output matched one candidate semantics among rivals; the replay checks
  consistency and is not independent evidence;
\item \emph{whole-program}: seen only inside programs verified bit-exact on the GPU, not run in isolation;
\item \emph{bounded negative}: searched for within a stated bound and not found; nothing is claimed beyond the bound;
\item \emph{Apple's compiler emits}: read from Apple's compiled code, which shows that a pattern works but leaves open
  whether the hardware requires it;
\item \emph{decoder}: what Apple's decoder prints or admits, with no execution behind it;
\item \emph{compile-witnessed}: a register number or form witnessed by compiling a source, with no dispatch confirming
  its values;
\item \emph{driver, static read}: read from the driver's code, not executed;
\item \emph{emulated}: run in the project's CPU emulator, which is checked against Apple's GPU;
\item \emph{this compiler}: a choice the project's compiler makes;
\item \emph{open}: unresolved.
\end{itemize}

Opcodes are named first and numbered second, as in "\op{12674}". The numbers are
Apple's opcode ids. The names are this project's, chosen by measured function, and are not Apple's spellings
(\cref{app:c}). Registers are written by hardware number (R0–R125), with their role stated where it matters.

Open questions are stated next to the rule they qualify and indexed in \cref{ch:not-known-index}. In Part~\ref{part:i}
a chapter with several open questions collects them in a paragraph headed \textbf{Not known}, and an \textbf{Evidence} list at the
end of a chapter, or of a long section, gathers the machine-model sections it cites.
